\documentclass[11pt,a4paper]{article}

\usepackage[margin=1in]{geometry}
\usepackage{amsmath,amssymb,amsthm}
\usepackage{mathtools}
\usepackage{booktabs}
\usepackage{algorithm}
\usepackage[noend]{algpseudocode}
\usepackage{hyperref}
\usepackage{xcolor}
\usepackage{microtype}

\hypersetup{colorlinks=true,linkcolor=blue!50!black,citecolor=blue!50!black,
            urlcolor=blue!50!black}

\theoremstyle{plain}
\newtheorem{theorem}{Theorem}[section]
\newtheorem{proposition}[theorem]{Proposition}
\newtheorem{lemma}[theorem]{Lemma}
\newtheorem{corollary}[theorem]{Corollary}
\theoremstyle{definition}
\newtheorem{definition}[theorem]{Definition}
\newtheorem{example}[theorem]{Example}
\newtheorem{algorithmx}[theorem]{Algorithm}
\theoremstyle{remark}
\newtheorem{remark}[theorem]{Remark}

\DeclareMathOperator{\spf}{spf}
\DeclareMathOperator{\rad}{rad}
\DeclareMathOperator{\artanh}{artanh}
\newcommand{\Z}{\mathbb{Z}}
\newcommand{\Q}{\mathbb{Q}}
\newcommand{\N}{\mathbb{N}}
\newcommand{\R}{\mathbb{R}}
\newcommand{\isqrt}{\operatorname{isqrt}}
\newcommand{\e}[1]{e\!\left(#1\right)}

\title{The Imbalance Sieve:\\
Enumeration, Novelty, and Spectral Structure of the Rationals}
\author{Paul A. Bilokon\thanks{Thalesians Ltd and Department of Mathematics,
Imperial College London. Email: \texttt{paul.bilokon@imperial.ac.uk}.}}
\date{August 2026}

\begin{document}
\maketitle

\begin{abstract}
For integers $1 \le p < q$ define the \emph{imbalance}
$\delta(p,q) = (q-p)/(q+p)$. Enumerating all such pairs by increasing $q$ and
then increasing $p$ produces a single sequence of rationals. We develop this
construction from first principles. The central object is the parity-adjusted
map $\Phi(x,y) = (A(y-x),\,A(y+x))$, where $A(n) = n/\gcd(n,2)$: we prove it is
an \emph{involution} on primitive pairs, and that consequently the passage
between a lattice point and a reduced fraction, in either direction, requires no
greatest common divisor. A value is \emph{novel} when it has not occurred
earlier, and this happens exactly when its generating pair is coprime; the novel
subsequence therefore enumerates $\Q \cap (0,1)$ bijectively, and is the Farey
enumeration in imbalance coordinates. We give exact algorithms with proved
complexity for every access pattern: positional access in $O(M(b))$, streaming
in $O(1)$ amortised, the novelty mask to index $N$ in $O(N)$ without a single
gcd, and ordinal access to the $n$-th novel term in $O(n^{1/3+\varepsilon})$.
The scalar sequence satisfies an explicit second-order recurrence, and two is
the minimal order. We show that the associated transition operator is the
unilateral shift and is therefore spectrally trivial, and that the genuine
spectral content lies in the emitted novelty mask: its Ramanujan expansion has
frequency support exactly the divisor lattice of $\rad(q)$, its autocorrelation
is a product over the primes dividing $q$, and its mean value is a
Hardy--Littlewood-type singular series with $r^2$ in place of $r$, hence without
parity obstruction. All results are accompanied by a verified reference
implementation.
\end{abstract}

\tableofcontents

\section{Introduction}

\subsection{The coordinate}

Given two positive quantities one may compare them by their difference, by their
ratio, or by the normalised difference
\begin{equation}\label{eq:imb}
  \delta(p,q) \;=\; \frac{q-p}{q+p}, \qquad 1 \le p < q .
\end{equation}
The third is scale-free: it depends only on $p/q$, and determines it, since
\begin{equation}\label{eq:inverse}
  \frac{p}{q} \;=\; \frac{1-\delta}{1+\delta}.
\end{equation}
Thus $\delta$ is a M\"obius transform of the ratio, and indeed an involution of
$(0,1)$: writing $\sigma(x) = (1-x)/(1+x)$ we have $\sigma \circ \sigma = \mathrm{id}$
and $\delta(p,q) = \sigma(p/q)$. It is also a logarithmic coordinate,
\begin{equation}\label{eq:artanh}
  \log \frac{q}{p} \;=\; \log\frac{1+\delta}{1-\delta} \;=\; 2\artanh \delta,
  \qquad |\delta| < 1,
\end{equation}
so that $\delta = \tanh\bigl(\tfrac12 \log(q/p)\bigr)$: the imbalance is the
hyperbolic tangent of half the logarithmic distance between $p$ and $q$. This
identity, exploited in Section~\ref{sec:metric}, is the reason the construction
has a metric interpretation at all.

The quantity is not exotic. In queueing and in market microstructure,
$(q-p)/(q+p)$ is the standard normalised imbalance of two competing quantities;
in relativity it is the rapidity coordinate; in elementary geometry, for
symmetric endpoints $p = n-r$, $q = n+r$ it is simply $r/n$, so that consecutive
odd integers $n-1, n+1$ have imbalance $1/n$.

\subsection{What this paper does}

We take the pairs $(p,q)$ with $1 \le p < q$, order them lexicographically by
$q$ and then $p$, and study the resulting rational sequence
$(\delta_i)_{i \ge 0}$. The questions are elementary and the answers are
uniformly explicit:

\begin{enumerate}
  \item \emph{Where is a given term?} Section~\ref{sec:index} inverts the
    triangular enumeration and gives an algorithm whose correctness rests on an
    integer comparison rather than on a square root, so that a floating-point
    seed cannot produce a wrong answer.
  \item \emph{How is a term reduced?} Section~\ref{sec:involution} isolates the
    parity adjustment $A$ and proves that $\Phi(x,y) = (A(y-x), A(y+x))$ is an
    involution on primitive pairs. Reduction and its inverse are the same map.
  \item \emph{When is a term new?} Section~\ref{sec:novelty} proves that a value
    is novel precisely when its pair is coprime, identifies the novel
    subsequence with the Farey enumeration, and records what this does and does
    not say about primes.
  \item \emph{Where does a given rational appear?} Section~\ref{sec:inverse}
    gives the full trail of occurrence indices of a fixed value as an explicit
    quadratic, and the fibres over a fixed numerator as interleaved quadratic
    progressions.
  \item \emph{Is the sequence self-generating?} Section~\ref{sec:recurrence}
    exhibits an explicit second-order recurrence and proves two is minimal.
  \item \emph{What is its spectrum?} Section~\ref{sec:spectral} shows the
    transition operator is spectrally trivial and computes the genuine spectral
    invariants of the emitted mask.
  \item \emph{Which rationals are born in prime frontiers?}
    Section~\ref{sec:partition} gives the criterion, the density, and the
    Dirichlet decomposition, and shows the induced partition is invisible to the
    natural limiting measure.
  \item \emph{What metric does this induce on the primes?}
    Section~\ref{sec:metric}.
\end{enumerate}

Section~\ref{sec:algorithms} collects the algorithms and their complexities,
Section~\ref{sec:numerics} the numerical verification, and
Section~\ref{sec:limits} an explicit account of what the construction does
\emph{not} deliver.

\subsection{Relation to classical enumerations}
\label{sec:classical}

The rationals have been enumerated many times: by Farey sequences
\cite{Farey1816}, by the mediant constructions of Stern \cite{Stern1858} and
Brocot \cite{Brocot1861}, and by the tree of Calkin and Wilf \cite{CalkinWilf2000}.
The present enumeration is not new as a set-theoretic device, and we state its
relation to the classical picture plainly in Proposition~\ref{prop:farey}: the
novel terms of frontier at most $N$ correspond bijectively to the Farey
fractions of order $N$ in $(0,1)$. What the imbalance coordinate contributes is
not a new enumeration but an explicit, gcd-free, cost-symmetric \emph{addressing
scheme} for it, together with the spectral invariants of
Section~\ref{sec:spectral}.

An earlier note by the author \cite{Bilokon2025} treated the unit-fraction case
of Section~\ref{sec:inverse}. We take the opportunity to correct an error in it;
see Remark~\ref{rem:erratum}.

Background on elementary divisibility and the totient is in
\cite{Apostol1976,HardyWright2008}; on integer arithmetic and square roots,
\cite{BrentZimmermann2021}; on Ramanujan expansions,
\cite{Ramanujan1918,Schwarz1994}; on sieve-theoretic context,
\cite{IwaniecKowalski2004}.

\section{The triangular enumeration}
\label{sec:index}

\subsection{Frontiers}

\begin{definition}[Frontier]
For $q \ge 2$, the \emph{$q$-th frontier} is the finite ordered list of pairs
$(1,q), (2,q), \dots, (q-1,q)$, and the corresponding list of imbalances
\[
  I(q) \;=\; \left( \frac{q-1}{q+1},\ \frac{q-2}{q+2},\ \dots,\ \frac{1}{2q-1} \right).
\]
\end{definition}

We order all pairs by frontier and, within a frontier, by increasing $p$:
\begin{equation}\label{eq:order}
  (1,2),\ (1,3),\ (2,3),\ (1,4),\ (2,4),\ (3,4),\ (1,5),\ \dots
\end{equation}
and write $(p_i,q_i)$ for the pair at index $i \ge 0$ and $\delta_i$ for its
imbalance. Let $T_m = m(m+1)/2$.

\begin{theorem}[Explicit index formula]\label{thm:index}
For $i \ge 0$ put
\begin{equation}\label{eq:mj}
  m = m(i) = \left\lfloor \frac{\sqrt{8i+1}-1}{2} \right\rfloor,
  \qquad j = j(i) = i - T_m .
\end{equation}
Then $0 \le j \le m$, and
\begin{equation}\label{eq:pair}
  (p_i, q_i) = (j+1,\ m+2), \qquad
  \delta_i = \frac{m+1-j}{m+3+j}.
\end{equation}
\end{theorem}

\begin{proof}
Frontier $q$ contains $q-1$ pairs, so it occupies the index block
$T_{q-2} \le i < T_{q-1}$. Hence $q_i = m+2$ where $m$ is determined by
$T_m \le i < T_{m+1}$, and $p_i = j+1$ with $j = i - T_m$, giving
$0 \le j \le m$. The condition $T_m \le i < T_{m+1}$ is
$m(m+1) \le 2i < (m+1)(m+2)$, equivalent to
$m \le (\sqrt{8i+1}-1)/2 < m+1$, which is~\eqref{eq:mj}. Substituting
$p = j+1$, $q = m+2$ in~\eqref{eq:imb} gives~\eqref{eq:pair}.
\end{proof}

\begin{table}[t]
\centering
\caption{The first $28$ imbalances, in lowest terms, grouped by frontier.
Repeated values are marked with a dagger; all others are first arrivals.}
\label{tab:first}
\begin{tabular}{cll}
\toprule
$q$ & indices & imbalance values \\
\midrule
$2$ & $0$ & $\tfrac13$ \\
$3$ & $1$--$2$ & $\tfrac12,\ \tfrac15$ \\
$4$ & $3$--$5$ & $\tfrac35,\ \tfrac13^\dagger,\ \tfrac17$ \\
$5$ & $6$--$9$ & $\tfrac23,\ \tfrac37,\ \tfrac14,\ \tfrac19$ \\
$6$ & $10$--$14$ & $\tfrac57,\ \tfrac12^\dagger,\ \tfrac13^\dagger,\ \tfrac15^\dagger,\ \tfrac1{11}$ \\
$7$ & $15$--$20$ & $\tfrac34,\ \tfrac59,\ \tfrac25,\ \tfrac3{11},\ \tfrac16,\ \tfrac1{13}$ \\
$8$ & $21$--$27$ & $\tfrac79,\ \tfrac35^\dagger,\ \tfrac5{11},\ \tfrac13^\dagger,\ \tfrac3{13},\ \tfrac17^\dagger,\ \tfrac1{15}$ \\
\bottomrule
\end{tabular}
\end{table}

\subsection{Exact positional access}

Formula~\eqref{eq:mj} contains a square root, and a naive implementation
computes $\lfloor \sqrt{8i+1} \rfloor$ exactly and then takes a further floor.
This is the wrong place to demand exactness: the composite floor is fragile
precisely at the triangular indices, where $8i+1 = (2m+1)^2$ is a perfect
square, and those are exactly the indices at which each frontier begins. The
following formulation moves the exactness requirement onto an integer
comparison, so that the square root becomes a mere seed.

\begin{theorem}[Seeded access with an integer certificate]\label{thm:seed}
Let $i \ge 0$ and let $m$ be as in~\eqref{eq:mj}. Then
\begin{equation}\label{eq:twocand}
  m \in \bigl\{\, s-1,\ s \,\bigr\}, \qquad s := \isqrt(2i) = \lfloor \sqrt{2i} \rfloor,
\end{equation}
and $m = s$ if $T_s \le i$, and $m = s-1$ otherwise. Moreover, for
\emph{any} integer seed $m_0 \ge 0$, the procedure
\[
  \textbf{while } m_0 > 0 \textbf{ and } T_{m_0} > i:\ m_0 \gets m_0 - 1;
  \qquad
  \textbf{while } T_{m_0+1} \le i:\ m_0 \gets m_0 + 1
\]
terminates after exactly $|m_0 - m|$ steps and returns $m$.
\end{theorem}

\begin{proof}
From $T_m \le i$ we get $m^2 \le m(m+1) \le 2i$, hence $m \le s$. From
$i < T_{m+1}$ we get $2i < (m+1)(m+2) < (m+2)^2$, hence $s < m+2$, i.e.
$s \le m+1$. Together, $m \in \{s-1, s\}$, and the case is decided by whether
$T_s \le i$. For the second claim, the map $m \mapsto T_m$ is strictly
increasing on $\N$, so the predicate $T_{m_0} \le i$ is monotone in $m_0$; each
loop moves $m_0$ one step towards the unique $m$ satisfying
$T_m \le i < T_{m+1}$, and neither loop can overshoot it.
\end{proof}

\begin{remark}
Theorem~\ref{thm:seed} separates correctness from cost. Correctness follows from
two multiplications and a comparison and holds for any seed whatsoever; the seed
quality only affects the number of correction steps. In particular a
double-precision seed $m_0 = \lfloor(\sqrt{8i+1}-1)/2\rfloor$ computed in
floating point is safe at any magnitude, and is within one of $m$ whenever the
conversion of $8i+1$ is exact, i.e. for $i < 2^{50}$.
\end{remark}

\begin{remark}[Comparison with a prescribed Babylonian count]
One may instead compute $\lfloor \sqrt{N} \rfloor$, $N = 8i+1$, by the integer
Babylonian iteration $B_N(x) = \min\{x, \lfloor (x + \lfloor N/x \rfloor)/2\rfloor\}$
started at $x_0 = N$; a standard error-halving argument shows
$\lceil \log_2 N\rceil + 1$ iterations suffice. Each iteration contains a
full-width division, so this costs $O(b\,M(b))$ with $b = \log N$, a factor $b$
more than necessary. Seeding at $x_0 = 2^{\lceil b/2\rceil}$ instead reduces the
count to $\lceil \log_2 b \rceil + 2$ by quadratic convergence, and the recursive
algorithm of \cite[Alg.~1.12]{BrentZimmermann2021} attains $O(M(b))$ outright.
Theorem~\ref{thm:seed} sidesteps the question by not requiring an exact root at
all.
\end{remark}

\begin{corollary}[Streaming]\label{cor:stream}
Define $F(p,q) = (p+1, q)$ if $p < q-1$ and $F(q-1,q) = (1, q+1)$. Then
$(p_{i+1}, q_{i+1}) = F(p_i, q_i)$, so consecutive terms are generated in $O(1)$
integer operations each, with no square root and no division.
\end{corollary}

\begin{proof}
Immediate from the ordering~\eqref{eq:order}: within frontier $q$ the index
advances $p$ by one, and at $p = q-1$ the frontier is exhausted and the next
begins at $(1,q+1)$.
\end{proof}

\section{The parity involution}
\label{sec:involution}

\subsection{The adjustment \texorpdfstring{$A$}{A}}

\begin{definition}\label{def:A}
Let
\[
  A(n) \;=\; \frac{n}{\gcd(n,2)} \;=\;
  \begin{cases} n, & n \text{ odd},\\ n/2, & n \text{ even}.\end{cases}
\]
This is OEIS \texttt{A026741} \cite{OEIS_A026741}.
\end{definition}

\begin{proposition}[Lowest-term formula]\label{prop:lowest}
Let $1 \le p < q$ and $h = \gcd(p,q)$. Then the lowest-term numerator and
denominator of $\delta(p,q)$ are
\begin{equation}\label{eq:lowest}
  \operatorname{num} \delta(p,q) = A\!\left(\frac{q-p}{h}\right), \qquad
  \operatorname{den} \delta(p,q) = A\!\left(\frac{q+p}{h}\right).
\end{equation}
\end{proposition}

\begin{proof}
Write $p = ha$, $q = hb$ with $\gcd(a,b) = 1$. Then
$\delta(p,q) = (b-a)/(b+a)$. Any common divisor $c$ of $b-a$ and $b+a$ divides
their sum $2b$ and their difference $2a$, hence divides $\gcd(2a,2b) = 2$. So
$c \in \{1,2\}$, and $c = 2$ occurs exactly when $b-a$ and $b+a$ are both even,
i.e. when $a$ and $b$ are both odd. Applying $A$ to each of $b-a$ and $b+a$
removes precisely this common factor of two when it is present and does nothing
otherwise.
\end{proof}

\subsection{The involution}

Call a pair $(x,y)$ with $0 < x < y$ \emph{primitive} if $\gcd(x,y) = 1$.

\begin{theorem}[$\Phi$ is an involution]\label{thm:involution}
Define $\Phi(x,y) = \bigl(A(y-x),\, A(y+x)\bigr)$. Then $\Phi$ maps primitive
pairs to primitive pairs, and $\Phi \circ \Phi = \mathrm{id}$.
\end{theorem}

\begin{proof}
Let $(x,y)$ be primitive and set $(u,v) = \Phi(x,y)$. Two cases.

\emph{Case 1: $x, y$ both odd.} Then $y-x$ and $y+x$ are even and
$(u,v) = \bigl(\tfrac{y-x}{2}, \tfrac{y+x}{2}\bigr)$, so $v-u = x$ and
$v+u = y$. Any common divisor of $u$ and $v$ divides $u+v = y$ and $v-u = x$,
hence divides $\gcd(x,y) = 1$; and $0 < u < v$ since $0 < x < y$. So $(u,v)$ is
primitive. Moreover $u + v = y$ is odd, so $u,v$ have opposite parity, and
therefore
\[
  \Phi(u,v) = \bigl(A(v-u), A(v+u)\bigr) = \bigl(A(x), A(y)\bigr) = (x,y),
\]
because $x$ and $y$ are odd.

\emph{Case 2: $x, y$ of opposite parity.} Then $y-x$ and $y+x$ are both odd and
$(u,v) = (y-x, y+x)$. Any common divisor $c$ of $u,v$ is odd and divides
$u+v = 2y$ and $v-u = 2x$, hence divides $\gcd(x,y) = 1$. Again $0<u<v$, so
$(u,v)$ is primitive, and both are odd, so
\[
  \Phi(u,v) = \bigl(A(v-u), A(v+u)\bigr) = \bigl(A(2x), A(2y)\bigr) = (x,y). \qedhere
\]
\end{proof}

Theorem~\ref{thm:involution} is the structural core of the paper. It says that
the passage from a primitive lattice point to a reduced fraction and the passage
back are \emph{the same map}, and that this map costs one subtraction, one
addition, one parity test and one shift. Everything gcd-free in what follows is
a consequence.

\begin{corollary}[Free reduction at coprime pairs]\label{cor:free}
If $\gcd(p,q) = 1$ then the lowest-term form of $\delta(p,q)$ is $\Phi(p,q)$,
computable in $O(b)$ bit operations, $b = \log q$, with no gcd. If
$\gcd(p,q) = h > 1$ then the scale $h$ itself is required
by~\eqref{eq:lowest}, and testing $h = 1$ does not suffice.
\end{corollary}

\begin{example}
At index $i = 12$ the pair is $(3,6)$, so $h = 3$ and the reduced value is
$1/3$; but $\Phi(3,6) = (A(3), A(9)) = (3,9)$, which is not reduced. Applying
$\Phi$ without checking coprimality is therefore incorrect in general. At
$i = 4$, with pair $(2,4)$ and $h = 2$, it happens to be correct, because $A$
removes exactly the factor $2$ that $h$ would have removed.
\end{example}

\section{Novelty}
\label{sec:novelty}

\begin{definition}[Novelty]
The term $\delta_i$ is \emph{novel} if $\delta_i \notin \{\delta_0, \dots,
\delta_{i-1}\}$. Write $\nu_i = \mathbf{1}\{\delta_i \text{ novel}\}$ for the
\emph{novelty mask}.
\end{definition}

\begin{theorem}[Novelty criterion]\label{thm:novelty}
For every $i \ge 0$,
\[
  \nu_i = 1 \iff \gcd(p_i, q_i) = 1 .
\]
Equivalently, the set of pairs decomposes as a disjoint union of dilation rays
$(p,q) = k\,(P,Q)$ with $(P,Q)$ primitive and $k \ge 1$, and only $k = 1$ is
novel.
\end{theorem}

\begin{proof}
By~\eqref{eq:inverse}, $\delta(p_1,q_1) = \delta(p_2,q_2)$ if and only if
$p_1/q_1 = p_2/q_2$. Each ratio in $(0,1)$ has a unique primitive representative
$(P,Q)$, and the pairs with that ratio are exactly $k(P,Q)$, $k \ge 1$. Since
the ordering~\eqref{eq:order} is by increasing upper coordinate and $Q < kQ$ for
$k > 1$, the primitive pair occurs strictly before all its dilates. Hence a term
is a first occurrence exactly when its pair is primitive.
\end{proof}

\begin{corollary}[Bijection with the rationals]\label{cor:bijection}
The novel subsequence enumerates $\Q \cap (0,1)$, each value exactly once. The
full sequence enumerates each value countably often, along its dilation ray.
\end{corollary}

\begin{proof}
By Theorem~\ref{thm:novelty} the novel terms correspond bijectively to primitive
pairs $(P,Q)$, $0 < P < Q$; by Theorem~\ref{thm:involution} these correspond
bijectively to reduced fractions in $(0,1)$, via $\Phi$.
\end{proof}

\begin{remark}[Erratum to \cite{Bilokon2025}]\label{rem:erratum}
Section~6.2 of \cite{Bilokon2025} asserts that
$\Phi(p,q) = (p-q)/(p+q)$ is injective on all pairs $p > q > 0$. It is not: the
proof correctly derives $p_1 q_2 = p_2 q_1$ and hence proportionality, but then
asserts the proportionality constant is $1$, which does not follow. The
counterexample $\delta(2,1) = \delta(4,2) = 1/3$ appears in that paper's own
Appendix~A. The correct statement is injectivity on \emph{primitive} pairs,
which is Theorem~\ref{thm:novelty} above, and it is what the surjectivity
argument of that section needs.
\end{remark}

\subsection{Relation to the Farey enumeration}

\begin{proposition}[The novel terms are Farey]\label{prop:farey}
For $N \ge 2$, the novel terms with frontier at most $N$ are in bijection, via
$\Phi$, with the Farey fractions of order $N$ lying in the open interval
$(0,1)$. Consequently their number is
\[
  \sum_{q=2}^{N} \varphi(q) \;=\; \Phi_\Sigma(N) - 1,
  \qquad \Phi_\Sigma(N) := \sum_{q \le N} \varphi(q) ,
\]
and the natural density of novel indices is
\[
  \lim_{N \to \infty} \frac{\Phi_\Sigma(N)-1}{T_{N-1}}
  \;=\; \frac{6}{\pi^2}.
\]
\end{proposition}

\begin{proof}
Novel terms with frontier at most $N$ correspond to primitive pairs
$0 < p < q \le N$, i.e. to reduced fractions $p/q \in (0,1)$ with $q \le N$,
which is the definition of $F_N \cap (0,1)$. The count is
$\sum_{q\le N}\varphi(q)$ less the term $\varphi(1) = 1$. The density follows
from $\Phi_\Sigma(N) = 3N^2/\pi^2 + O(N\log N)$ and $T_{N-1} \sim N^2/2$.
\end{proof}

So the construction is the Farey enumeration in imbalance coordinates. What is
added is the addressing: Section~\ref{sec:inverse} gives the index of a
prescribed fraction in closed form, which the mediant constructions do not.

\subsection{Arithmetic read-outs}

\begin{corollary}[Blockwise arithmetic]\label{cor:arith}
For every $q \ge 2$,
\begin{equation}\label{eq:phi}
  \varphi(q) \;=\; \sum_{i=T_{q-2}}^{T_{q-1}-1} \nu_i,
\end{equation}
and
\begin{equation}\label{eq:prime}
  q \text{ is prime} \iff \nu_i = 1 \text{ for all } T_{q-2} \le i < T_{q-1}.
\end{equation}
If $q$ is composite then
\begin{equation}\label{eq:spf}
  \spf(q) \;=\; \min\{\,p \ge 2 \;:\; \nu_{T_{q-2}+p-1} = 0 \,\}.
\end{equation}
\end{corollary}

\begin{proof}
The frontier $q$ consists of the pairs $(p,q)$, $1 \le p < q$, and by
Theorem~\ref{thm:novelty} the novel ones are those with $\gcd(p,q) = 1$; there
are $\varphi(q)$ of them, giving~\eqref{eq:phi}. All are novel exactly when
every $p < q$ is coprime to $q$, i.e. $q$ is prime, giving~\eqref{eq:prime}. If
$q$ is composite, the least $p \ge 2$ with $\gcd(p,q) > 1$ is the least prime
factor of $q$, giving~\eqref{eq:spf}.
\end{proof}

\begin{remark}[What Corollary~\ref{cor:arith} does and does not say]
\label{rem:honest}
Each of \eqref{eq:phi}--\eqref{eq:spf} is an exact identity, and iterating
\eqref{eq:spf} recovers the full factorisation of every frontier label. It is
important to be clear that this is a change of representation and not a
computational advance: $\nu$ is \emph{defined} through $\gcd$, so possessing the
mask is not cheaper than possessing the factorisations it encodes. What is
genuine is the cost statement in the other direction --- the mask can be
generated from a smallest-prime-factor sieve at $O(1)$ amortised cost per entry,
with no gcd at all (Algorithm~\ref{alg:mask} and Theorem~\ref{thm:masklinear}),
so the two representations are interconvertible at known cost. Similarly, with
$C(q) = \prod_{i=T_{q-2}}^{T_{q-1}-1}\nu_i$ we have the exact equivalence
\[
  q,\, q+2 \text{ both prime} \iff C(q)\,C(q+2) = 1,
\]
but this is a restatement of the twin prime conjecture in different notation,
and confers no leverage on it.
\end{remark}

\subsection{Imbalance trails}

The frontier admits a description in terms of numerators which makes the
appearance of $A$ transparent.

\begin{definition}[Imbalance trail]
For $q \ge 2$ the \emph{trail} of $q$ is the sequence $\tau_k(q) = \delta(q-k,q)$
for $k = 1, \dots, q-1$, in lowest terms, with reduced numerator $A_k(q)$.
\end{definition}

\begin{proposition}[Prime trail pattern]\label{prop:trail}
Let $q \ge 5$. Then $q$ is prime if and only if
\begin{equation}\label{eq:trail}
  A_k(q) = A(k) =
  \begin{cases} k, & k \text{ odd},\\ k/2, & k \text{ even},\end{cases}
  \qquad 1 \le k \le q-1 .
\end{equation}
\end{proposition}

\begin{proof}
Here $p = q - k$, so $q - p = k$ and $q + p = 2q - k$, and
$h = \gcd(q-k, q) = \gcd(k,q)$. By Proposition~\ref{prop:lowest},
$A_k(q) = A(k/h)$. If $q$ is prime then for $1 \le k \le q-1$ we have
$q \nmid k$, so $h = 1$ and $A_k(q) = A(k)$, which is~\eqref{eq:trail}.
Conversely, if $q$ is composite let $r = \spf(q) \le q-1$; then $h = \gcd(r,q) = r$
and $A_r(q) = A(1) = 1 \ne A(r)$ since $r \ge 2$. Hence~\eqref{eq:trail} fails
at $k = r$.
\end{proof}

Thus the numerator sequence of a prime trail is the initial segment of
\texttt{A026741}, and the first deviation of a composite trail occurs at
$k = \spf(q)$. This is~\eqref{eq:prime} and~\eqref{eq:spf} read numerator-wise;
it is the same content, and it explains why $A$ is the right adjustment to have
introduced in Definition~\ref{def:A}.

\begin{remark}
Checking~\eqref{eq:trail} for all $k$ is trial division: the condition
$\gcd(k,q) \in \{1,2\}$ for all $k < q$ is exactly the statement that no
$k < q$ shares an odd factor with $q$, and the first violation is at
$k = \spf(q)$. The trail formulation is a reinterpretation of trial division, of
the same worst-case cost $\widetilde{O}(\sqrt q)$ and the same
polylogarithmic average cost, not a new primality test.
\end{remark}

\section{Inverse indexing}
\label{sec:inverse}

\subsection{Occurrence trails}

\begin{proposition}[Occurrence indices]\label{prop:occ}
Let $a/b \in (0,1)$ be reduced, and set $(P,Q) = \Phi(a,b) = (A(b-a), A(b+a))$.
Then $(P,Q)$ is the unique primitive pair with $\delta(P,Q) = a/b$, the pairs
producing $a/b$ are exactly $(kP, kQ)$ for $k \ge 1$, and the $k$-th occurrence
has index
\begin{equation}\label{eq:occ}
  i_k \;=\; T_{kQ-2} + kP - 1 \;=\; \frac{k^2 Q^2 + k(2P - 3Q)}{2}.
\end{equation}
In particular the unique novel occurrence is
\begin{equation}\label{eq:novidx}
  i_{\mathrm{nov}} \;=\; \frac{Q^2 + 2P - 3Q}{2}.
\end{equation}
\end{proposition}

\begin{proof}
Solving $(q-p)/(q+p) = a/b$ gives $p/q = (b-a)/(b+a)$. By the argument of
Proposition~\ref{prop:lowest} the only possible common factor of $b-a$ and
$b+a$ is $2$, removed exactly by $A$; hence $(P,Q) = \Phi(a,b)$ is primitive and
generates $a/b$, and it is unique by Theorem~\ref{thm:involution}. All integer
solutions are its dilates. The pair $(kP,kQ)$ lies at index
$T_{kQ-2} + kP - 1$ by Theorem~\ref{thm:index}; expanding
$T_{kQ-2} = (k^2Q^2 - 3kQ + 2)/2$ gives~\eqref{eq:occ}. Novelty at $k=1$ is
Theorem~\ref{thm:novelty}.
\end{proof}

\begin{corollary}[Constant second difference]\label{cor:seconddiff}
The occurrence indices of a fixed reduced $a/b$ satisfy
\[
  i_{k+2} - 2 i_{k+1} + i_k \;=\; Q^2 ,
\]
and $\#\{k : i_k \le x\} = \sqrt{2x}/Q + O(1)$.
\end{corollary}

\begin{proof}
Both follow from~\eqref{eq:occ}, a quadratic in $k$ with leading coefficient
$Q^2/2$.
\end{proof}

\begin{example}\label{ex:trails}
The first three occurrence trails are
\[
\begin{aligned}
  \mathrm{Occ}(1/3) &: \ i_k = 2k(k-1) = 0, 4, 12, 24, 40, \dots \\
  \mathrm{Occ}(1/2) &: \ i_k = \tfrac{9k^2 - 7k}{2} = 1, 11, 30, 58, 95, \dots \\
  \mathrm{Occ}(3/5) &: \ i_k = 8k^2 - 5k = 3, 22, 57, 108, 175, \dots
\end{aligned}
\]
with second differences $4, 9, 16$, the squares of the primitive upper
coordinates $2, 3, 4$.
\end{example}

\subsection{Fixed numerator}

Fixing the numerator and letting the denominator vary gives a finite family of
quadratic progressions.

\begin{theorem}[Interleaved progressions]\label{thm:progressions}
Fix $a \ge 1$ and let $b$ range over integers $b > a$ with $\gcd(a,b) = 1$. Put
$M = 2\rad(a)$. Then the residue of $b$ modulo $M$ takes exactly
$2\varphi(\rad(a))$ admissible values, and on each such residue class the novel
index of $a/b$ is a quadratic polynomial in $b$, namely
\begin{equation}\label{eq:prog}
  i_{\mathrm{nov}}(a,b) =
  \begin{dcases}
    \frac{(b+a)^2 + 4(b-a) - 6(b+a)}{8}, & a, b \text{ both odd},\\[4pt]
    \frac{(b+a)^2 + 2(b-a) - 3(b+a)}{2}, & a, b \text{ of opposite parity}.
  \end{dcases}
\end{equation}
\end{theorem}

\begin{proof}
Coprimality of $a$ and $b$ depends only on $b \bmod \rad(a)$, and holds for
$\varphi(\rad(a))$ residues modulo $\rad(a)$; refining modulo $M = 2\rad(a)$
doubles this to $2\varphi(\rad(a))$ classes, on each of which the parity of $b$
is constant. By Proposition~\ref{prop:occ}, $i_{\mathrm{nov}} = (Q^2 + 2P - 3Q)/2$
with $(P,Q) = ((b-a)/\varepsilon, (b+a)/\varepsilon)$ and $\varepsilon = 2$ when
$a,b$ are both odd and $\varepsilon = 1$ otherwise. Substituting gives the two
displayed expressions, each a quadratic in $b$.
\end{proof}

\begin{corollary}[Unit fractions]\label{cor:unit}
For $a = 1$ there are exactly two classes, and
\[
  i_{\mathrm{nov}}(1,d) =
  \begin{cases} \dfrac{d^2-9}{8}, & d \text{ odd},\\[6pt]
                \dfrac{d(d+1)}{2} - 2, & d \text{ even}.\end{cases}
\]
Moreover $1/d$ is the first novel imbalance with reduced denominator $d$, with
primitive pair $(P_d,Q_d) = (d-1,d+1)$ for $d$ even and
$\bigl(\tfrac{d-1}{2},\tfrac{d+1}{2}\bigr)$ for $d$ odd.
\end{corollary}

\begin{proof}
$\rad(1) = 1$, so $2\varphi(\rad(1)) = 2$ and the two classes are the parities of
$d$; substituting $a=1$ in~\eqref{eq:prog} gives the displayed formulas. For the
second claim, a reduced $a/d$ has primitive upper coordinate $Q = A(d+a)$ by
Proposition~\ref{prop:occ}, and by~\eqref{eq:occ} its novel index is increasing
in $Q$ for fixed parity class. If $d$ is even then $a$ is odd and
$A(d+a) = d+a$, minimised uniquely at $a=1$. If $d$ is odd and $a$ is odd then
$A(d+a) = (d+a)/2$, again minimised at $a=1$, and any even $a$ gives the larger
value $d+a \ge d+2 > (d+1)/2$. Substituting $a=1$ yields the stated pairs.
\end{proof}

Corollary~\ref{cor:unit} is Theorem 5.1 of \cite{Bilokon2025}, recovered here as
the case $a = 1$ of Theorem~\ref{thm:progressions}.

\section{A second-order recurrence}
\label{sec:recurrence}

The sequence $(\delta_i)$ is generated by the deterministic map $F$ of
Corollary~\ref{cor:stream} acting on the pair $(p,q)$. That pair is, however,
merely the index in disguise: Theorem~\ref{thm:index} is a bijection between $i$
and $(p_i, q_i)$. A more informative statement is that the scalar sequence
regenerates itself from its own past.

\begin{theorem}[Order two, and no less]\label{thm:order2}
For $i \ge 1$ the pair $(p_i, q_i)$ is determined by $\delta_{i-1}$ and
$\delta_i$ alone:
\begin{equation}\label{eq:recover}
  (p_i, q_i) =
  \begin{dcases}
    \left( q\,\frac{1-\delta_i}{1+\delta_i},\ q \right),\quad
      q = \frac{(1+\delta_i)(1+\delta_{i-1})}{2(\delta_{i-1}-\delta_i)},
      & \delta_i < \delta_{i-1},\\[6pt]
    \left( 1,\ \frac{1+\delta_i}{1-\delta_i} \right), & \delta_i > \delta_{i-1}.
  \end{dcases}
\end{equation}
Consequently there is an explicit $G : \Q^2 \to \Q$ with
$\delta_{i+1} = G(\delta_{i-1}, \delta_i)$ for all $i \ge 1$. No first-order
recurrence exists: $\delta$ alone is not a state.
\end{theorem}

\begin{proof}
Write $x_i = p_i/q_i = \sigma(\delta_i)$. Within a frontier, $p$ increases by
one at fixed $q$, so $x_i - x_{i-1} = 1/q_i > 0$ and hence
$\delta_i < \delta_{i-1}$, since $\sigma$ is strictly decreasing. At a frontier
boundary, $\delta$ jumps from $1/(2q-1)$ to $q/(q+2)$, and
$q/(q+2) > 1/(2q-1)$ for $q \ge 2$, so $\delta_i > \delta_{i-1}$. The two cases
are therefore distinguished by the sign of $\delta_i - \delta_{i-1}$, which
never vanishes.

In the first case,
\[
  x_i - x_{i-1}
  = \frac{1-\delta_i}{1+\delta_i} - \frac{1-\delta_{i-1}}{1+\delta_{i-1}}
  = \frac{2(\delta_{i-1}-\delta_i)}{(1+\delta_i)(1+\delta_{i-1})}
  = \frac{1}{q_i},
\]
which inverts to the stated $q$, and then $p_i = q_i x_i$. In the second case
$p_i = 1$ and $\delta_i = (q_i-1)/(q_i+1)$ inverts to $q_i = (1+\delta_i)/(1-\delta_i)$.
Applying $F$ to the recovered pair and taking the imbalance gives $G$.

For minimality, $\delta_0 = \delta_4 = 1/3$ while $\delta_1 = 1/2$ and
$\delta_5 = 1/7$: equal terms with unequal successors, so no function of
$\delta_i$ alone can determine $\delta_{i+1}$.
\end{proof}

\begin{remark}
The analogy is with the Fibonacci recurrence, where the augmented state
$(F_n, F_{n+1})$ is built from the sequence's own outputs. Theorem~\ref{thm:order2}
is the correct analogue: the scale information apparently destroyed by reduction
is recovered from one additional observation. This also sharpens the applied
reading of $\delta$ as a normalised queue imbalance: two consecutive readings
determine the underlying integer pair, whereas one does not.
\end{remark}

\section{Spectral structure}
\label{sec:spectral}

\subsection{Why the transition operator is uninformative}

\begin{proposition}[Degeneracy of the transition spectrum]\label{prop:degenerate}
Let $X_i = (p_i,q_i)$ with $X_{i+1} = F(X_i)$. Then $(X_i)$ is a
time-homogeneous deterministic Markov chain on a countable state space; its
transition operator $(\mathcal{P}f)(x) = f(F(x))$ acting on $\ell^2$ is
unitarily equivalent to the \emph{backward} unilateral shift $S^*$.
Consequently $\sigma(\mathcal{P}) = \overline{\mathbb{D}}$, every
$\lambda$ in the open unit disc is an eigenvalue of $\mathcal{P}$, the adjoint
$\mathcal{P}^* = S$ has no eigenvalues at all, and the chain admits no non-zero
invariant measure. The same holds for the induced chains on the novel
subsequence and on the novel indices.
\end{proposition}

\begin{proof}
By Theorem~\ref{thm:index} the state space is in bijection with $\N$ under
$i \leftrightarrow (p_i,q_i)$, and under this bijection $F$ is $i \mapsto i+1$;
hence $(\mathcal{P}f)_i = f_{i+1}$, which is $S^*$. The equation
$S^* f = \lambda f$ gives $f_i = \lambda^i f_0$, which lies in $\ell^2$ for
every $|\lambda| < 1$, so the point spectrum of $\mathcal{P}$ is the open disc;
the spectrum is closed and contained in the unit ball, hence equals
$\overline{\mathbb{D}}$. For $S$, the equation $Sf = \lambda f$ forces
$f_0 = 0$ and then $f_i = 0$ for all $i$, so $S$ has no eigenvalues. Finally
$F$ is injective with $(1,2)$ outside its image, so $\mu = \mu \circ F^{-1}$
forces $\mu(\{(1,2)\}) = 0$, and then inductively $\mu(\{F^n(1,2)\}) = 0$ for
every $n$; on a countable state space this gives $\mu \equiv 0$. Restricting to
a subsequence replaces $F$ by a first-return map, which is again an injective
shift on a countable set.
\end{proof}

The degeneracy is thus as complete as it could be: the spectrum is the largest
it can possibly be, and every point of its interior is an eigenvalue, so no
spectral quantity distinguishes this sequence from any other.

Every deterministic sequence, over any alphabet, has this transition spectrum:
take the state space $\N$, the map $i \mapsto i+1$, and the emission
$i \mapsto a_i$. No ergodic, mixing or spectral-gap statement is available or
meaningful here. The spectral content of the construction lies entirely in the
\emph{emissions}, and is almost periodic rather than dynamical. We compute it
now.

\subsection{Ramanujan expansion of the novelty mask}

Fix a frontier $q$ and regard $\nu$ as the function
$\nu(p) = \mathbf{1}\{\gcd(p,q) = 1\}$ on $\Z/q$. Write
$\e{t} = \exp(2\pi i t)$ and let
$c_d(k) = \sum_{\substack{1\le t \le d,\ \gcd(t,d)=1}} \e{tk/d}$
be the Ramanujan sum, so that $c_d(k) = \mu(d/g)\varphi(d)/\varphi(d/g)$ with
$g = \gcd(k,d)$.

\begin{theorem}[Ramanujan expansion]\label{thm:ramanujan}
For every $q \ge 2$ and every $p$,
\begin{equation}\label{eq:ram}
  \nu(p) \;=\; \frac{\varphi(q)}{q} \sum_{d \mid q} \frac{\mu(d)}{\varphi(d)}\, c_d(p).
\end{equation}
\end{theorem}

\begin{proof}
The finite Fourier coefficients of $\nu$ on $\Z/q$ are
\[
  \widehat\nu(k) = \frac1q \sum_{p \bmod q} \nu(p)\, \e{-pk/q}
  = \frac1q \sum_{\gcd(p,q)=1} \e{-pk/q} = \frac{c_q(k)}{q},
\]
using that $c_q$ is real and even. Hence
$\nu(p) = q^{-1}\sum_{k \bmod q} c_q(k) \e{pk/q}$. Group the $k$ by
$g = \gcd(k,q)$ and put $d = q/g$. For fixed $g$, writing $k = g k'$ with
$\gcd(k', d) = 1$,
\[
  \sum_{\gcd(k,q) = g} \e{pk/q} = \sum_{\substack{k' \bmod d \\ \gcd(k',d)=1}} \e{pk'/d} = c_d(p),
\]
while $c_q(k) = \mu(d)\varphi(q)/\varphi(d)$ is constant on that class. Summing
over $d \mid q$ gives~\eqref{eq:ram}.
\end{proof}

\begin{corollary}[The spectrum sees only the radical]\label{cor:radical}
The frequency support of~\eqref{eq:ram} is
$\{ d : d \mid \rad(q) \}$, of size $2^{\omega(q)}$, and the amplitudes
$\varphi(q)\mu(d)/(q\varphi(d))$ depend only on $\rad(q)$, since
$\varphi(q)/q = \prod_{r \mid q}(1 - 1/r)$. In particular the masks of $q$ and
of $\rad(q)$ have identical spectra. Moreover, given the frontier label $q$
(which is determined by the frontier length),
\[
  q \text{ is prime} \iff \text{the support is } \{1, q\}.
\]
\end{corollary}

\begin{proof}
$\mu(d) = 0$ unless $d$ is squarefree, and the squarefree divisors of $q$ are
the divisors of $\rad(q)$. The stated dependence of the amplitudes is immediate.
The divisors of $\rad(q)$ are $\{1,q\}$ exactly when $q$ is prime: if $q = r^e$
with $e \ge 2$ they are $\{1,r\} \ne \{1,q\}$, and if $q$ has at least two
distinct prime factors there are at least four.
\end{proof}

Corollary~\ref{cor:radical} explains why~\eqref{eq:spf} requires a recursion:
multiplicities are invisible to a single frontier's spectrum.

\subsection{Autocorrelation}

\begin{theorem}[Exact autocorrelation]\label{thm:autocorr}
For $q \ge 2$ and $h \in \Z$ let
$R_q(h) = q^{-1} \#\{ p \bmod q : \gcd(p,q) = \gcd(p+h,q) = 1 \}$. Then
\begin{equation}\label{eq:autocorr}
  R_q(h) \;=\; \prod_{r \mid \gcd(h,q)} \left(1 - \frac1r\right)
               \prod_{\substack{r \mid q \\ r \nmid h}} \left(1 - \frac2r\right),
\end{equation}
the products being over primes.
\end{theorem}

\begin{proof}
By the Chinese remainder theorem the count factorises over the prime powers
$r^e \,\|\, q$. Modulo $r^e$, the condition is $r \nmid p$ and $r \nmid p+h$. If
$r \mid h$ these coincide and the count is $r^e(1-1/r)$. If $r \nmid h$ they
exclude two distinct residues mod $r$, giving $r^e(1-2/r)$. Multiplying and
dividing by $q$ gives~\eqref{eq:autocorr}.
\end{proof}

\begin{corollary}
$R_q(0) = \varphi(q)/q$; and if $q$ is even and $h$ is odd then $R_q(h) = 0$,
since the factor at $r=2$ vanishes.
\end{corollary}

\subsection{The singular series}

\begin{theorem}[Mean value]\label{thm:singular}
For fixed $h$, the multiplicative function $q \mapsto R_q(h)$ has mean value
\begin{equation}\label{eq:sing}
  \mathfrak{S}(h) \;=\; \prod_{r \mid h}\left(1 - \frac{1}{r^2}\right)
                        \prod_{r \nmid h}\left(1 - \frac{2}{r^2}\right),
\end{equation}
and consequently
\[
  \sum_{q \le Q} q\,R_q(h) \;\sim\; \mathfrak{S}(h)\,\frac{Q^2}{2},
  \qquad Q \to \infty .
\]
In particular $\mathfrak{S}(0) = \prod_r (1 - r^{-2}) = 6/\pi^2$, since every
prime divides $0$.
\end{theorem}

\begin{proof}
Write $f(q) = R_q(h) = \prod_{r \mid q} \alpha_r$ with $\alpha_r = 1 - 1/r$ if
$r \mid h$ and $\alpha_r = 1 - 2/r$ otherwise; by Theorem~\ref{thm:autocorr}
this is multiplicative, bounded by $1$, and satisfies $f(r^e) = \alpha_r$ for
all $e \ge 1$.

Put $g = \mu * f$, so that $f = \mathbf{1} * g$. Then $g$ is multiplicative with
$g(r) = \alpha_r - 1$ and $g(r^e) = f(r^e) - f(r^{e-1}) = 0$ for $e \ge 2$;
hence $g$ is supported on squarefree integers, with
$g(d) = \prod_{r \mid d}(\alpha_r - 1)$ there. Since
$|\alpha_r - 1| \le 2/r$, we get $|g(d)| \le 2^{\omega(d)}/d$, so
\begin{equation}\label{eq:gabs}
  \sum_{d} \frac{|g(d)|}{d} \le \prod_r \left(1 + \frac{2}{r^2}\right) < \infty,
  \qquad \sum_{d \le D} |g(d)| \le \sum_{d\le D} \frac{2^{\omega(d)}}{d} = O(\log^2 D).
\end{equation}
Therefore
\[
  \sum_{q \le Q} q\, f(q)
  = \sum_{d \le Q} g(d)\, d \sum_{m \le Q/d} m
  = \frac{Q^2}{2} \sum_{d \le Q} \frac{g(d)}{d}
    + O\!\left( Q \sum_{d\le Q} |g(d)| \right)
  = \frac{Q^2}{2} \sum_{d} \frac{g(d)}{d} + O(Q \log^2 Q),
\]
the tail being controlled by~\eqref{eq:gabs}. Finally, by multiplicativity,
\[
  \sum_{d} \frac{g(d)}{d}
  = \prod_r \left(1 + \frac{\alpha_r - 1}{r}\right)
  = \prod_r \left(1 - \frac1r + \frac{\alpha_r}{r}\right),
\]
and $1 - 1/r + \alpha_r/r$ equals $1 - r^{-2}$ when $r \mid h$ and $1 - 2r^{-2}$
otherwise, which is~\eqref{eq:sing}.
\end{proof}

\begin{remark}[No parity obstruction]\label{rem:parity}
Expression~\eqref{eq:sing} has the shape of a Hardy--Littlewood singular series
with $r^2$ in place of $r$. The consequence is structural: the factor at $r = 2$
is $1 - 2/4 = 1/2 \ne 0$ for odd $h$, so no factor of $\mathfrak{S}$ vanishes and
the correlations are strictly positive for every lag. Even lags exceed odd lags
by exactly the ratio
$\mathfrak{S}(2)/\mathfrak{S}(1) = (1-\tfrac14)/(1-\tfrac12) = \tfrac32$. The
novelty mask therefore has genuine arithmetic memory, increasing with the small
prime divisors of the lag.
\end{remark}

\section{The prime-born partition}
\label{sec:partition}

Each rational in $(0,1)$ is born in exactly one frontier
(Corollary~\ref{cor:bijection}). We may therefore partition $\Q \cap (0,1)$ by
the primality of that frontier.

\begin{definition}
A reduced $a/b \in (0,1)$ is \emph{prime-born} if its birth frontier
$Q = A(a+b)$ is prime, and \emph{composite-born} otherwise.
\end{definition}

\begin{proposition}[Criterion and density]\label{prop:primeborn}
$a/b$ is prime-born if and only if $A(a+b)$ is prime. The proportion of
prime-born values among those born in frontiers up to $x$ is
\[
  \frac{\sum_{q \le x,\ q \text{ prime}} (q-1)}{\sum_{q\le x} \varphi(q)}
  \;\sim\; \frac{\pi^2}{6 \log x},
\]
so the prime-born set has natural density zero in this ordering.
\end{proposition}

\begin{proof}
The criterion is Proposition~\ref{prop:occ}. A prime frontier $q$ contributes
all of its $q-1$ entries, since $\varphi(q) = q-1$. By the prime number theorem
and partial summation $\sum_{p \le x} (p-1) \sim x^2/(2\log x)$, while
$\sum_{q\le x}\varphi(q) \sim 3x^2/\pi^2$. The ratio is
$(\pi^2/6)/\log x$.
\end{proof}

Primes give \emph{complete} frontiers but are too sparse to carry positive
density: completeness buys the constant $\pi^2/6$, sparsity costs a factor
$\log x$.

\begin{proposition}[Dirichlet decomposition]\label{prop:dirichlet}
Weighting each novel value by its birth frontier, for $\Re s > 2$,
\[
  \sum_{\delta \text{ novel}} Q(\delta)^{-s}
  = \sum_{q \ge 2} \varphi(q)\, q^{-s}
  = \frac{\zeta(s-1)}{\zeta(s)} - 1 ,
\]
and this splits as
\[
  \underbrace{P(s-1) - P(s)}_{\text{prime-born}}
  \;+\;
  \underbrace{\frac{\zeta(s-1)}{\zeta(s)} - 1 - P(s-1) + P(s)}_{\text{composite-born}},
\]
where $P$ is the prime zeta function. At $s = 2$ the total has a simple pole
with residue $6/\pi^2$, while the prime-born part has only a logarithmic
singularity; this is Proposition~\ref{prop:primeborn} read off from the order of
the singularity.
\end{proposition}

\begin{proof}
$\sum_{q\ge1}\varphi(q)q^{-s} = \zeta(s-1)/\zeta(s)$ is classical; subtract the
term $q=1$. The prime-born values are those with $Q$ prime, contributing
$\sum_{r} (r-1) r^{-s} = P(s-1) - P(s)$.
\end{proof}

\begin{proposition}[The partition is invisible to the limiting measure]
\label{prop:measures}
Let $\mu_x$ be the uniform probability measure on the novel imbalances born in
frontiers at most $x$, and $\mu_x^{\mathrm{pr}}$ the same for the prime-born
ones. Then both converge weakly, as $x \to \infty$, to the measure on $(0,1)$
with density
\[
  \frac{2}{(1+\delta)^2}\,d\delta .
\]
\end{proposition}

\begin{proof}
The Farey fractions of order $x$ equidistribute with respect to Lebesgue measure
on $(0,1)$; for prime $q$ the frontier contributes the complete grid
$\{p/q : 1 \le p \le q-1\}$, which also equidistributes, and prime frontiers are
infinite in number. In both cases the ratios $p/q$ equidistribute uniformly, and
the imbalances are their images under $\sigma(x) = (1-x)/(1+x)$. If $X$ is
uniform on $(0,1)$ and $Y = \sigma(X)$ then $X = \sigma(Y)$ and
$|dX/dY| = 2/(1+Y)^2$, which integrates to $1$ over $(0,1)$.
\end{proof}

Thus the two classes are distributionally identical and differ only in total
mass. What \emph{is} visible at the level of a single frontier is finer, and is
the correct statement of \eqref{eq:prime}: frontier $q$ is prime exactly when
its measure is the complete uniform grid of spacing $1/q$, and composite
frontiers are that grid punctured at the multiples of the prime divisors of $q$.
Primality is minimal local discrepancy.

\section{The induced metric on the primes}
\label{sec:metric}

\begin{proposition}[$\delta$ is a metric]\label{prop:metric}
For $x,y > 0$ set $d(x,y) = |y-x|/(y+x)$, with $d(x,x)=0$. Then $d$ is a metric
on $\R_{>0}$, and
\[
  d(x,y) = \tanh\!\left( \tfrac12 \left| \log \tfrac{y}{x} \right| \right).
\]
\end{proposition}

\begin{proof}
The identity is~\eqref{eq:artanh}. Symmetry and positivity are clear. For the
triangle inequality, $\tanh$ is concave and increasing on $[0,\infty)$ with
$\tanh 0 = 0$, hence subadditive: $\tanh(u+v) \le \tanh u + \tanh v$. Writing
$\ell(x,y) = \tfrac12|\log(y/x)|$, which is a metric, and using monotonicity,
\[
  d(x,z) = \tanh \ell(x,z) \le \tanh\bigl(\ell(x,y)+\ell(y,z)\bigr)
  \le \tanh\ell(x,y) + \tanh\ell(y,z) = d(x,y) + d(y,z). \qedhere
\]
\end{proof}

Restrict $d$ to the set $\mathbb{P}$ of primes.

\begin{proposition}[Structure of $(\mathbb{P},d)$]\label{prop:primespace}
The space $(\mathbb{P}, d)$ is bounded of diameter $1$, the diameter not being
attained; it is discrete and complete; it is not totally bounded; and it is not
uniformly discrete, since
\[
  d(p_n, p_{n+1}) = \frac{g_n}{2p_n + g_n} \longrightarrow 0,
  \qquad g_n = p_{n+1}-p_n .
\]
\end{proposition}

\begin{proof}
$d < 1$ always and $d(2, p) \to 1$ as $p \to \infty$, giving the diameter. A
$d$-Cauchy sequence has $\log p_n$ Cauchy in $\R$, hence convergent, hence
eventually constant since $\{\log p\}$ is discrete and closed in $\R$; so the
space is discrete and complete. It is not totally bounded: $d(x,y)<\varepsilon$
forces $\log(y/x)$ into a bounded interval, and $\{\log p\}$ is unbounded, so no
finite family of $\varepsilon$-balls covers $\mathbb{P}$ for $\varepsilon < 1$.
The gap formula is direct.
\end{proof}

\begin{proposition}[Twins, and the arithmetic of distances]\label{prop:twins}
For primes $p < q$:
\begin{enumerate}
  \item $d(p,q)$ is always a novel imbalance, since $\gcd(p,q)=1$;
  \item $q = p+2$ if and only if $d(p,q) = 1/(p+1)$;
  \item if $p,q$ are odd, $d(p,q)$ is prime-born if and only if $(p+q)/2$ is
    prime.
\end{enumerate}
\end{proposition}

\begin{proof}
(i) is Theorem~\ref{thm:novelty}. (ii) $d(p,p+2) = 2/(2p+2) = 1/(p+1)$, and
conversely $d$ determines the ratio $q/p$ and hence, with $p$ fixed, $q$. (iii)
The birth frontier of $d(p,q)$ is $A(q+p)$, and for odd $p,q$ this is
$(p+q)/2$; apply Proposition~\ref{prop:primeborn}.
\end{proof}

Part (iii) says the distance between two odd primes is prime-born exactly when
their average is prime, i.e. exactly when $p+q = 2n$ with $n$ itself prime. This
is a statement about which Goldbach representations have prime half-sum; it is
not the Goldbach conjecture, and we claim nothing more for it.

\section{Algorithms and complexity}
\label{sec:algorithms}

Throughout, $b$ denotes the bit length of the index $i$, and $M(b)$ the cost of
$b$-bit multiplication.

\begin{algorithm}[t]
\caption{Positional access: the $i$-th entry}
\label{alg:entry}
\begin{algorithmic}[1]
\Require $i \ge 0$
\State $s \gets \isqrt(2i)$ \Comment{seed; any seed is admissible}
\If{$s(s+1)/2 \le i$} $m \gets s$ \Else{} $m \gets s-1$ \EndIf
  \Comment{integer certificate, Thm.~\ref{thm:seed}}
\State $j \gets i - m(m+1)/2$
\State \Return $(p,q) \gets (j+1,\, m+2)$ and $\delta \gets (m+1-j)/(m+3+j)$
\end{algorithmic}
\end{algorithm}

\begin{algorithm}[t]
\caption{The novelty mask on $[0,N)$, without gcd}
\label{alg:mask}
\begin{algorithmic}[1]
\Require $N \ge 1$
\State $Q \gets$ frontier of index $N-1$ \Comment{Algorithm~\ref{alg:entry}}
\State sieve smallest prime factors up to $Q$ \Comment{linear sieve}
\State $\nu_i \gets 1$ for all $i < N$
\For{$q = 2, \dots, Q$}
  \State $\text{base} \gets T_{q-2}$; factor $q$ from the spf table
  \ForAll{primes $r \mid q$}
    \For{$p = r, 2r, 3r, \dots < q$}
      \State $\nu_{\text{base}+p-1} \gets 0$
    \EndFor
  \EndFor
\EndFor
\end{algorithmic}
\end{algorithm}

\begin{theorem}[Linear-time mask]\label{thm:masklinear}
Algorithm~\ref{alg:mask} is correct and runs in $O(N)$ word operations and
$O(\sqrt N)$ space, invoking no gcd.
\end{theorem}

\begin{proof}
Correctness is Theorem~\ref{thm:novelty}: within frontier $q$, the non-novel
positions are exactly those $p$ divisible by some prime factor of $q$. The
frontier labels reach $Q \asymp \sqrt{2N}$, so the sieve costs $O(\sqrt N)$ time
and space. The number of marks is
\[
  \sum_{q \le Q}\ \sum_{r \mid q,\ r \text{ prime}} \frac{q}{r}
  \;=\; \sum_{r \le Q} \frac1r \sum_{\substack{q \le Q\\ r \mid q}} q
  \;\sim\; \frac{Q^2}{2}\sum_{r} \frac{1}{r^2}
  \;\approx\; 0.4522\,\frac{Q^2}{2},
\]
using $\sum_r r^{-2} = P(2) \approx 0.4522$, against $T_{Q-1} \sim Q^2/2$
indices. Hence $O(N)$ marks in total.
\end{proof}

\begin{algorithm}[t]
\caption{Ordinal access: the $n$-th novel index}
\label{alg:ordinal}
\begin{algorithmic}[1]
\Require $n \ge 1$
\State $q \gets \lceil \pi\sqrt{n/3}\,\rceil$
  \Comment{seed from $\Phi_\Sigma(Q) \sim 3Q^2/\pi^2$}
\While{$\Phi_\Sigma(q) - 1 < n$} $q \gets q+1$ \EndWhile
\While{$q > 2$ \textbf{and} $\Phi_\Sigma(q-1) - 1 \ge n$} $q \gets q-1$ \EndWhile
\State $r \gets n - (\Phi_\Sigma(q-1) - 1)$ \Comment{rank within frontier $q$}
\State $p \gets$ the $r$-th integer in $[1,q)$ coprime to $q$
  \Comment{binary search on $\sum_{d \mid \rad(q)} \mu(d)\lfloor x/d\rfloor$}
\State \Return $T_{q-2} + p - 1$
\end{algorithmic}
\end{algorithm}

\begin{theorem}[Cost of ordinal access]\label{thm:ordinal}
Algorithm~\ref{alg:ordinal} returns the $n$-th novel index in
$O(n^{1/3+\varepsilon})$ time and $O(n^{1/4})$ space.
\end{theorem}

\begin{proof}
By Proposition~\ref{prop:farey} the novel entries in frontiers up to $q$ number
$\Phi_\Sigma(q)-1$, so the loop terminates at the correct frontier. Evaluating
$\Phi_\Sigma(x)$ by the hyperbola recurrence
$\sum_{k \le x} \Phi_\Sigma(\lfloor x/k\rfloor) = T_x$, memoised over the
$O(\sqrt x)$ distinct values of $\lfloor x/k \rfloor$, costs $O(x^{2/3})$ time
and $O(\sqrt x)$ space. Here $q \asymp \sqrt n$, giving $O(n^{1/3})$ and
$O(n^{1/4})$. The totative search evaluates a sum over the $2^{\omega(q)}
= q^{o(1)}$ squarefree divisors of $q$, $O(\log q)$ times.
\end{proof}

\begin{table}[t]
\centering
\caption{Cost of each access pattern. Only ordinal access is genuinely
arithmetic; the others are combinatorial and certify in two multiplications.}
\label{tab:costs}
\begin{tabular}{llll}
\toprule
Operation & Method & Cost & gcd? \\
\midrule
$i \mapsto (p,q)$, unreduced $\delta_i$ & Alg.~\ref{alg:entry} & $O(M(b))$ & no \\
streaming & Cor.~\ref{cor:stream} & $O(1)$ amortised & no \\
lowest terms, $i$ novel & $\Phi$, Thm.~\ref{thm:involution} & $O(b)$ & no \\
lowest terms, streaming & $\Phi$ + per-frontier factorisation & $O(\omega(q))$ divisions & no Euclid \\
lowest terms, arbitrary $i$ & Prop.~\ref{prop:lowest} & $O(M(b)\log b)$ & yes \\
reduced $a/b \mapsto$ novel index & Prop.~\ref{prop:occ} & $O(M(b))$ & no \\
novelty mask on $[0,N)$ & Alg.~\ref{alg:mask} & $O(N)$, $O(\sqrt N)$ space & no \\
$n \mapsto n$-th novel index & Alg.~\ref{alg:ordinal} & $O(n^{1/3+\varepsilon})$ & no \\
\bottomrule
\end{tabular}
\end{table}

Table~\ref{tab:costs} exhibits the asymmetry that we regard as the operational
content of the construction. Addressing a novel term by \emph{name} --- its
rational value --- is cheap in both directions and symmetric, dominated by a
square root one way and a squaring the other. Addressing it by \emph{position}
costs a totient summatory computation. The bijection of
Corollary~\ref{cor:bijection} is explicit; its ordering is not.

\section{Numerical verification}
\label{sec:numerics}

A reference implementation accompanies this paper; all statements below were
checked against it.

\paragraph{Indexing and trails.} Algorithm~\ref{alg:entry} agrees with the
streaming recurrence of Corollary~\ref{cor:stream} for all $i < 2\cdot10^5$, and
reproduces Table~\ref{tab:first} entry by entry. The occurrence trails of
Example~\ref{ex:trails} and their second differences $4,9,16$ are confirmed, as
are the two progressions of Corollary~\ref{cor:unit} and the four progressions
predicted by Theorem~\ref{thm:progressions} for $a = 3$.

\paragraph{Mask and ordinal access.} Algorithm~\ref{alg:mask} agrees with the
gcd definition on $[0, 3\cdot 10^5)$, with empirical novelty density
$0.60991$ against $6/\pi^2 = 0.60793$. Algorithm~\ref{alg:ordinal} agrees with
the sieve for all sampled ranks up to $n = 1.2\cdot10^6$; at
$n = 10^{15}$ it returns $1\,644\,934\,015\,772\,573$ in under half a second,
against the prediction $n\pi^2/6 = 1.6449\cdots\times10^{15}$.

\paragraph{Spectral quantities.} Expansion~\eqref{eq:ram} reproduces $\nu$
exactly for all $q < 60$; the support claim of Corollary~\ref{cor:radical} is
confirmed, including that $q = 8$ and $q = 2$ have identical spectra, and the
primality criterion holds for all $q < 50$. Formula~\eqref{eq:autocorr} agrees
with direct counting for all $q < 80$ and all lags.

\paragraph{The singular series.} Table~\ref{tab:singular} compares
$\bigl(\sum_{q\le Q} q R_q(h)\bigr)/(Q^2/2)$ at $Q = 20\,000$ with
$\mathfrak{S}(h)$ of Theorem~\ref{thm:singular}.

\begin{table}[h]
\centering
\caption{Empirical mean of $R_q(h)$ at $Q = 20\,000$ against $\mathfrak{S}(h)$.}
\label{tab:singular}
\begin{tabular}{rlll}
\toprule
$h$ & empirical & $\mathfrak{S}(h)$ & ratio \\
\midrule
$0$  & $0.607952$ & $0.607927$ & $1.00004$ \\
$1$  & $0.322656$ & $0.322634$ & $1.00007$ \\
$2$  & $0.483979$ & $0.483952$ & $1.00006$ \\
$3$  & $0.368735$ & $0.368725$ & $1.00003$ \\
$6$  & $0.553105$ & $0.553087$ & $1.00003$ \\
$30$ & $0.577153$ & $0.577135$ & $1.00003$ \\
\bottomrule
\end{tabular}
\end{table}

The row $h=0$ returns $6/\pi^2$ as required, and
$\mathfrak{S}(2)/\mathfrak{S}(1) = 1.4999$, confirming
Remark~\ref{rem:parity}.

\paragraph{The prime-born partition.} Convergence in
Proposition~\ref{prop:primeborn} is slow, being of order $1/\log x$, and not
monotone at small scales:

\begin{center}
\begin{tabular}{rrrrr}
\toprule
$x$ & novel & prime-born & fraction & $\pi^2/(6\log x)$ \\
\midrule
$50$  & $773$   & $313$   & $0.4049$ & $0.4205$ \\
$100$ & $3043$  & $1035$  & $0.3401$ & $0.3572$ \\
$200$ & $12231$ & $4181$  & $0.3418$ & $0.3105$ \\
$400$ & $48677$ & $13809$ & $0.2837$ & $0.2745$ \\
\bottomrule
\end{tabular}
\end{center}

\section{Limitations}
\label{sec:limits}

We state plainly what the construction does not achieve.

\begin{enumerate}
  \item \emph{No new arithmetic.} Identities \eqref{eq:phi}--\eqref{eq:spf} and
    the twin-prime equivalence of Remark~\ref{rem:honest} are exact
    reformulations. The novelty mask is defined through $\gcd$, so it is not a
    cheaper route to factorisation; Theorem~\ref{thm:masklinear} shows only that
    the two representations are interconvertible in linear time.
  \item \emph{The prime-born partition is isomorphic to the prime/composite
    partition of $\N$}, by Corollary~\ref{cor:bijection}. Every statement about
    it is a translation, and Proposition~\ref{prop:measures} shows it is
    invisible to the natural limiting measure.
  \item \emph{Ordinal access admits no closed form.} By
    Proposition~\ref{prop:farey}, locating the $n$-th novel index requires
    $\Phi_\Sigma$ exactly, whose error term is
    $\Omega_\pm(x \log\log\log x)$ by Montgomery \cite{Montgomery1987} and is
    not known in order. Theorem~\ref{thm:ordinal} is therefore essentially the
    best available, and no radical expression in $n$ can exist. This is a
    genuine obstruction, not a gap in the present treatment.
  \item \emph{The dynamical framing is empty.} Proposition~\ref{prop:degenerate}
    holds for every deterministic sequence whatsoever. We retain the chain
    language only because Theorem~\ref{thm:order2} gives it content at the level
    of the scalar sequence.
  \item \emph{The metric of Section~\ref{sec:metric} is $\log$ in disguise},
    being $\tanh$ of half the logarithmic distance. Only
    Proposition~\ref{prop:twins}(iii), which ties prime-bornness of a distance to
    the primality of a half-sum, uses the arithmetic of the imbalance rather
    than the topology of the logarithm.
\end{enumerate}

\section{Open questions}

\begin{enumerate}
  \item Determine the second-order term in
    Proposition~\ref{prop:primeborn}, i.e. the size of the fluctuation visible
    in the table of Section~\ref{sec:numerics}.
  \item Is the exponent $1/3$ in Theorem~\ref{thm:ordinal} optimal for ordinal
    access, given the best known bounds for $\Phi_\Sigma$?
  \item The correlation $\mathfrak{S}(h)$ of Theorem~\ref{thm:singular} is a
    singular series without parity obstruction. Does the corresponding point
    process of novel indices have a limiting pair-correlation function, and is it
    the one predicted by $\mathfrak{S}$?
  \item In the applied reading of $\delta$ as a normalised queue imbalance with
    bounded integer queues, Theorem~\ref{thm:novelty} says the observable
    determines the pair only up to scale, and Theorem~\ref{thm:order2} says two
    consecutive observations suffice. Quantify the resulting quantisation:
    which imbalance values are attainable at depth $N$, with what multiplicities,
    and what is the induced discretisation error?
\end{enumerate}

\section*{Acknowledgements}

The author thanks Michel Marcus and Yasushi Ieno for corrections to the earlier
note \cite{Bilokon2025}.

\begin{thebibliography}{99}

\bibitem{Apostol1976} T. M. Apostol.
  \emph{Introduction to Analytic Number Theory}. Springer, New York, 1976.

\bibitem{Bilokon2025} P. A. Bilokon.
  A sieve on rational imbalances and the first appearance of denominators.
  SSRN preprint 5662230, 2025.

\bibitem{BrentZimmermann2021} R. P. Brent and P. Zimmermann.
  \emph{Modern Computer Arithmetic}. Cambridge University Press, 2nd ed., 2021.

\bibitem{Brocot1861} A. Brocot.
  Calcul des rouages par approximation, nouvelle m\'ethode.
  \emph{Revue Chronom\'etrique}, 3:186--194, 1861.

\bibitem{CalkinWilf2000} N. Calkin and H. S. Wilf.
  Recounting the rationals. \emph{Amer. Math. Monthly}, 107(4):360--363, 2000.

\bibitem{GrahamKnuthPatashnik1994} R. L. Graham, D. E. Knuth and O. Patashnik.
  \emph{Concrete Mathematics}. Addison-Wesley, 2nd ed., 1994.

\bibitem{HardyWright2008} G. H. Hardy and E. M. Wright.
  \emph{An Introduction to the Theory of Numbers}. Oxford University Press,
  6th ed., 2008.

\bibitem{IwaniecKowalski2004} H. Iwaniec and E. Kowalski.
  \emph{Analytic Number Theory}. AMS Colloquium Publications 53, 2004.

\bibitem{Farey1816} J. Farey.
  On a curious property of vulgar fractions.
  \emph{Philosophical Magazine}, 47:385--386, 1816.

\bibitem{Montgomery1987} H. L. Montgomery.
  Fluctuations in the mean of Euler's phi function.
  \emph{Proc. Indian Acad. Sci. Math. Sci.}, 97:239--245, 1987.

\bibitem{OEIS_A026741} OEIS Foundation Inc.
  Sequence A026741, \emph{The On-Line Encyclopedia of Integer Sequences},
  \url{https://oeis.org/A026741}.

\bibitem{Ramanujan1918} S. Ramanujan.
  On certain trigonometrical sums and their applications in the theory of
  numbers. \emph{Trans. Cambridge Philos. Soc.}, 22:259--276, 1918.

\bibitem{Schwarz1994} W. Schwarz and J. Spilker.
  \emph{Arithmetical Functions}. London Math. Soc. Lecture Notes 184,
  Cambridge University Press, 1994.

\bibitem{Stern1858} M. A. Stern.
  \"Uber eine zahlentheoretische Funktion.
  \emph{J. Reine Angew. Math.}, 55:193--220, 1858.

\end{thebibliography}

\end{document}
